
\section{Análisis y Especificación del Sistema}

En este capítulo se detallan las especificaciones completas del sistema MoodNest. Se definirán los actores involucrados, se detallarán los Requisitos Funcionales (RF) que determinan qué debe hacer el sistema, y los Requisitos No Funcionales (RNF) que establecen cómo debe comportarse técnica y operativamente.

\subsection{Análisis de Requisitos}

\subsubsection{Requisitos Funcionales (RF)}

Los Requisitos Funcionales describen el comportamiento esperado del sistema y las interacciones específicas que los usuarios pueden realizar dentro de la plataforma. Para una mejor comprensión, se han agrupado en módulos lógicos dentro de la siguiente tabla maestra.

\renewcommand{\arraystretch}{1.5}
\begin{xltabular}{\textwidth}{|>{\columncolor{white}\bfseries}c|X|}
    \hline
    \rowcolor{UCO}
    \textcolor{white}{\textbf{ID}} & \multicolumn{1}{c|}{\textcolor{white}{\textbf{Descripción del Requisito}}} \\ \hline
    \endfirsthead

    \multicolumn{2}{c}{{\bfseries \dots continuación de la página anterior}} \\
    \hline
    \rowcolor{UCO}
    \textcolor{white}{\textbf{ID}} & \multicolumn{1}{c|}{\textcolor{white}{\textbf{Descripción del Requisito}}} \\ \hline
    \endhead

    % --- MÓDULO 1 ---
    \multicolumn{2}{|c|}{\cellcolor{gray!20}\textbf{Módulo 1: Gestión de Usuarios}} \\ \hline
    RF1.1 & El sistema permitirá a un usuario no registrado crear una nueva cuenta proporcionando un nombre, un correo electrónico válido y único en el sistema, y una contraseña. \\ \hline
    RF1.2 & El sistema permitirá el acceso a usuarios registrados solicitando obligatoriamente su correo electrónico y contraseña. \\ \hline
    RF1.3 & El sistema permitirá al usuario autenticado modificar su nombre de visualización y actualizar su contraseña. \\ \hline
    RF1.4 & El sistema permitirá al usuario autenticado cerrar sesión de forma segura. \\ \hline
    RF1.5 & El sistema permitirá al usuario autenticado eliminar su cuenta y todos sus datos de forma permanente. \\ \hline

    % --- MÓDULO 2 ---
    \multicolumn{2}{|c|}{\cellcolor{gray!20}\textbf{Módulo 2: Registro Diario y Etiquetas}} \\ \hline
    RF2.1 & El sistema permitirá al usuario autenticado crear un nuevo registro de estado de ánimo asociado a una fecha. \\ \hline
    RF2.2 & El sistema establecerá la fecha actual por defecto al iniciar la creación de un nuevo registro, permitiendo su modificación únicamente a fechas pasadas y bloqueando explícitamente la selección de fechas futuras. \\ \hline
    RF2.3 & El sistema exigirá de forma obligatoria una puntuación global del estado de ánimo, representada por un valor numérico entero del 1 al 10. \\ \hline
    RF2.4 & El sistema permitirá al usuario incluir un comentario opcional en su registro diario. \\ \hline
    RF2.5 & El sistema permitirá al usuario autenticado crear nuevas etiquetas personalizadas asignándoles un nombre identificativo. \\ \hline
    RF2.6 & El sistema permitirá al usuario gestionar sus etiquetas personalizadas, pudiendo editarlas o eliminarlas de su catálogo. \\ \hline
    RF2.7 & El sistema permitirá al usuario asociar una o múltiples etiquetas a un registro diario. \\ \hline
    RF2.8 & El sistema permitirá al usuario asignar, de manera opcional, una puntuación numérica entera (del 1 al 10) a cada etiqueta individual seleccionada en un registro. \\ \hline
    RF2.9 & El sistema permitirá al usuario editar los atributos modificables (puntuación global, fecha asignada, comentario y etiquetas vinculadas) de un registro propio, aplicando las mismas reglas de validación y obligatoriedad definidas para la creación del registro. \\ \hline
    RF2.10 & El sistema permitirá al usuario eliminar de forma permanente cualquiera de sus registros pasados. \\ \hline

    % --- MÓDULO 3 ---
    \multicolumn{2}{|c|}{\cellcolor{gray!20}\textbf{Módulo 3: Panel Principal e Historial}} \\ \hline
    RF3.1 & El sistema mostrará al usuario autenticado un panel principal con el resumen de su último registro de estado de ánimo. \\ \hline
    RF3.2 & El sistema proporcionará en el panel principal un botón de acceso directo a la creación de un nuevo registro, accesible sin necesidad de navegación adicional. \\ \hline
    RF3.3 & El sistema mostrará en el panel principal un listado cronológico resumido con las últimas entradas registradas por el usuario. \\ \hline
    RF3.4 & El sistema calculará y mostrará visualmente el número de días consecutivos que el usuario ha realizado al menos un registro. \\ \hline
    RF3.5 & El sistema proporcionará al usuario autenticado una vista de historial basada en un calendario mensual, donde se identificarán visualmente los días que contienen un registro de estado de ánimo. \\ \hline
    RF3.6 & El sistema permitirá al usuario interactuar con el calendario, de forma que al seleccionar un día con registro, se muestre el detalle completo de la entrada. \\ \hline

    % --- MÓDULO 4 ---
    \multicolumn{2}{|c|}{\cellcolor{gray!20}\textbf{Módulo 4: Análisis de Datos (Estadísticas)}} \\ \hline
    RF4.1 & El sistema generará un gráfico temporal de líneas que represente la evolución de la puntuación global del estado de ánimo. \\ \hline
    RF4.2 & El sistema permitirá al usuario filtrar los datos de dicha evolución temporal seleccionando el rango de tiempo deseado (semanal, mensual o anual). \\ \hline
    RF4.3 & El sistema calculará y mostrará qué etiquetas están asociadas matemáticamente con las puntuaciones globales medias más altas (impacto positivo) y más bajas (impacto negativo) del estado de ánimo del usuario. \\ \hline
    RF4.4 & El sistema permitiría al usuario seleccionar una etiqueta específica de su historial para analizar su impacto individual. \\ \hline
    RF4.5 & El sistema generará un gráfico de barras agrupadas que compare visualmente el promedio del estado de ánimo global con el promedio de la puntuación específica de la etiqueta seleccionada. \\ \hline
    RF4.6 & El sistema generará automáticamente un mensaje de texto interpretativo basado en la diferencia matemática entre la puntuación global media y la puntuación de la etiqueta analizada. \\ \hline
    RF4.7 & El sistema generará un gráfico de barras que represente el promedio histórico de la puntuación global de los registros de estado de ánimo de cada uno de los días de la semana. \\ \hline
    RF4.8 & El sistema generará un gráfico circular que represente la distribución porcentual de los registros de estado de ánimo en diferentes rangos (ej. días buenos, neutros y malos). \\ \hline

    % --- MÓDULO 5 ---
    \multicolumn{2}{|c|}{\cellcolor{gray!20}\textbf{Módulo 5: Personalización y Sistema}} \\ \hline
    RF5.1 & El sistema permitirá al usuario autenticado personalizar los textos descriptivos de su escala de estado de ánimo para cada uno de los niveles numéricos (del 1 al 10). \\ \hline
    RF5.2 & El sistema permitirá al usuario seleccionar el estilo visual de su escala de ánimo eligiendo entre varios conjuntos de iconos predefinidos. \\ \hline
    RF5.3 & El sistema mostrará en el panel principal un mensaje de bienvenida dinámico determinado por dos parámetros: la franja horaria local y el estado de su registro diario (pendiente o completado). \\ \hline
    RF5.4 & El sistema permitirá al usuario alternar la interfaz gráfica entre un tema visual claro y un tema visual oscuro. \\ \hline
    RF5.5 & El sistema permitirá al usuario seleccionar un color principal para personalizar la apariencia de los elementos gráficos clave de la interfaz. \\ \hline
    \caption{Catálogo unificado de Requisitos Funcionales} \label{tab:rf_total} \\

\end{xltabular}

\newpage

\subsubsection{Requisitos No Funcionales (RNF)}

Los Requisitos No Funcionales definen los atributos de calidad del sistema, condicionando las decisiones tecnológicas del desarrollo de MoodNest.

\renewcommand{\arraystretch}{1.5}
\begin{xltabular}{\textwidth}{|>{\columncolor{white}\bfseries}c|X|}
    \hline
    \rowcolor{UCO}
    \textcolor{white}{\textbf{ID}} & \multicolumn{1}{c|}{\textcolor{white}{\textbf{Descripción del Requisito}}} \\ \hline
    \endfirsthead

    \multicolumn{2}{c}{{\bfseries \dots continuación de la página anterior}} \\
    \hline
    \rowcolor{UCO}
    \textcolor{white}{\textbf{ID}} & \multicolumn{1}{c|}{\textcolor{white}{\textbf{Descripción del Requisito}}} \\ \hline
    \endhead

    % --- BLOQUE 1 ---
    \multicolumn{2}{|c|}{\cellcolor{gray!20}\textbf{Bloque 1: Compatibilidad y Portabilidad}} \\ \hline
    RNF1.1 & El sistema deberá estar desarrollado bajo una arquitectura web Cliente-Servidor. \\ \hline
    RNF1.2 & El sistema deberá ser completamente accesible y funcional a través de los navegadores web modernos más utilizados. \\ \hline
    RNF1.3 & El sistema no requerirá la instalación de ningún software adicional ni plugins por parte del usuario. \\ \hline

    % --- BLOQUE 2 ---
    \multicolumn{2}{|c|}{\cellcolor{gray!20}\textbf{Bloque 2: Rendimiento y Disponibilidad}} \\ \hline
    RNF2.1 & El tiempo de respuesta del servidor para operaciones estándar no superará los 2 segundos. \\ \hline
    RNF2.2 & La interfaz gráfica principal deberá renderizarse en un tiempo máximo de 3 segundos. \\ \hline
    RNF2.3 & El servidor deberá ser capaz de procesar al menos 50 peticiones concurrentes por segundo. \\ \hline
    RNF2.4 & Ante falta de respuesta de la base de datos (5s), el sistema mostrará un mensaje de error amigable evitando el bloqueo. \\ \hline

    % --- BLOQUE 3 ---
    \multicolumn{2}{|c|}{\cellcolor{gray!20}\textbf{Bloque 3: Usabilidad e Interfaz}} \\ \hline
    RNF3.1 & La interfaz de usuario deberá ser Responsive (Móvil y Escritorio). \\ \hline
    RNF3.2 & El sistema permitirá acceder a la pantalla de creación de registro en un máximo de dos interacciones desde el inicio de sesión. \\ \hline
    RNF3.3 & El sistema proporcionará retroalimentación visual en un tiempo máximo de 0,5 segundos ante cualquier acción importante. \\ \hline

    % --- BLOQUE 4 ---
    \multicolumn{2}{|c|}{\cellcolor{gray!20}\textbf{Bloque 4: Seguridad y Privacidad}} \\ \hline
    RNF4.1 & El sistema almacenará las contraseñas aplicando un algoritmo de hashing unidireccional. \\ \hline
    RNF4.2 & El sistema protegerá las vistas privadas exigiendo un token de autenticación válido. \\ \hline
    RNF4.3 & El servidor verificará que el usuario solo acceda a sus propios registros y etiquetas. \\ \hline
    RNF4.4 & Toda la comunicación de datos se realizará de forma cifrada mediante el protocolo HTTPS. \\ \hline

    \caption{Catálogo unificado de Requisitos No Funcionales} \label{tab:rnf_total} \\

\end{xltabular}

\newpage

% --- COMANDO PARA TABLAS DE REQUISITOS ---

\newcommand{\reqTable}[3]{
\begin{table}[H]
\centering
\renewcommand{\arraystretch}{1.5} % Da un poco de "aire" a las filas
\begin{tabularx}{\textwidth}{|>{\columncolor{gray!20}\bfseries}c|X|}
\hline
\rowcolor{UCO} % Color de la cabecera (definido en confg_colores)
\multicolumn{2}{|c|}{\textcolor{white}{\textbf{#1}}} \\
\hline
#2
\end{tabularx}
\caption{#3}
\end{table}
}

% ------------------------------------------------



\subsubsection{Requisitos de Información (RI)}
\label{tab:ri_total}

Los Requisitos de Información definen el modelo conceptual de datos del sistema. Describen qué información debe almacenar, gestionar y recuperar la plataforma para garantizar su correcto funcionamiento, independientemente de la tecnología de base de datos que se seleccione posteriormente en la fase de diseño.

\reqTable{RI-01: Información del Usuario}{
    ID\_Usuario & Identificador único y autogenerado por el sistema. \\ \hline
    Nombre & Cadena de caracteres con el nombre o apodo visible en el panel. \\ \hline
    Email & Dirección de correo electrónico (debe ser un valor único en todo el sistema). \\ \hline
    Contraseña & Credencial de acceso almacenada de forma segura mediante un algoritmo de cifrado irreversible. \\ \hline
    Preferencias\_Sistema & Conjunto de configuraciones de interfaz, que incluye: el tema visual (claro/oscuro), el color principal y la familia de iconos seleccionada. \\ \hline
    Escala\_Personalizada & Conjunto de 10 descriptores textuales asociados a cada nivel numérico del estado de ánimo. \\ \hline
    Fecha\_Registro & Marca temporal (fecha y hora) exacta de la creación de la cuenta. \\ \hline
    Racha\_Actual & Valor numérico entero que contabiliza los días consecutivos realizando registros. \\ \hline
    Fecha\_Ultimo\_Registro & Fecha correspondiente a la última entrada creada en el sistema. \\ \hline
}{Requisitos de Información - Usuario}

\reqTable{RI-02: Información de la Etiqueta}{
    ID\_Etiqueta & Identificador único y autogenerado por el sistema. \\ \hline
    ID\_Usuario & Identificador que vincula la etiqueta con su usuario creador. \\ \hline
    Nombre & Cadena de caracteres con el nombre descriptivo de la actividad o factor (ej. "Deporte", "Familia"). \\ \hline
    Color & Código de color asociado a la etiqueta. \\ \hline
    Activa & Indicador lógico que determina si la etiqueta es visible y utilizable por el usuario, o si ha sido archivada/eliminada. \\ \hline
    Fecha\_Ultimo\_Uso & Marca temporal que se actualiza cada vez que la etiqueta se vincula a un nuevo registro, permitiendo ordenar la interfaz por uso reciente. \\ \hline
    Fecha\_Creacion & Marca temporal de la creación de la etiqueta. \\ \hline
}{Requisitos de Información - Etiqueta}

\reqTable{RI-03: Información del Registro Diario}{
    ID\_Registro & Identificador único y autogenerado por el sistema. \\ \hline
    ID\_Usuario & Identificador que vincula el registro con el usuario propietario. \\ \hline
    Fecha\_Asignada & Fecha del calendario correspondiente al día que se está evaluando. \\ \hline
    Puntuacion\_Global & Valor numérico entero (del 1 al 10) que representa el estado de ánimo general del día. \\ \hline
    Comentario & Texto opcional con anotaciones o descripciones adicionales del día. \\ \hline
    Etiquetas\_Asociadas & Relación de etiquetas vinculadas a este registro. Por cada etiqueta seleccionada, se guardará su identificador y, de forma opcional, una puntuación numérica específica (del 1 al 10). \\ \hline
    Fecha\_Creacion & Marca temporal exacta en la que el registro fue guardado en el sistema. \\ \hline
    Fecha\_Modificacion & Marca temporal de la última edición realizada sobre el registro. \\ \hline
}{Requisitos de Información - Registro Diario}

\newpage
\subsection{Verificación y Validación de Requisitos}

Una vez desarrollado por completo el catálogo de requisitos, llevamos a cabo un proceso de Verificación y Validación. El objetivo principal de esta etapa es asegurar la calidad, viabilidad y corrección del catálogo de requisitos antes de pasar a la siguiente fase. Aplicar esta técnica de forma temprana permite detectar problemas o incongruencias desde el principio, minimizando el impacto temporal y el coste que supondría arrastrar estos fallos a las fases de diseño o programación.

Esta validación se ha aplicado al catálogo completo de requisitos, incluyendo los Requisitos Funcionales (RF), los Requisitos No Funcionales (RNF) y los Requisitos de Información (RI). 

Para realizar la evaluación, se ha utilizado una lista de control (\textit{Checklist}) de calidad. Cada requisito se ha revisado individualmente para comprobar que cumple con estos cinco criterios básicos:

\begin{itemize}
    \item \textbf{C1 - Inequívoco:} El requisito está redactado con un lenguaje formal de forma clara y precisa. No usa términos subjetivos garantizando que admite una única interpretación técnica posible.
    \item \textbf{C2 - Verificable:} La especificación permite el diseño de una prueba objetiva que determine de forma concluyente si el sistema desarrollado cumple con la condición exigida.
    \item \textbf{C3 - Necesario:} El requisito aporta valor real a la aplicación y está alineado con los objetivos del proyecto, evitando añadir funciones de relleno o innecesarias.
    \item \textbf{C4 - Consistente:} El requisito no está repetido y no entra en conflicto ni se contradice con ningún otro punto del documento.
    \item \textbf{C5 - Completo:} La regla de negocio o restricción técnica está totalmente acotada e incluye toda la información necesaria para su desarrollo, definiendo correctamente sus límites sin delegar decisiones clave a fases posteriores.
\end{itemize}

% Comando para las celdas de la matriz
\newcommand{\pass}{\cmark} % Tick verde
\newcommand{\fail}{\xmark} % Cruz roja
\newcommand{\warn}{\textcolor{orange}{\textbf{!}}} % Advertencia naranja

\subsubsection{Matriz de Verificación: Requisitos Funcionales}

\renewcommand{\arraystretch}{1.3}
\begin{xltabular}{\textwidth}{|c|>{\centering\arraybackslash}X|>{\centering\arraybackslash}X|>{\centering\arraybackslash}X|>{\centering\arraybackslash}X|>{\centering\arraybackslash}X|c|}
    \hline
    \rowcolor{UCO}
    \textcolor{white}{\textbf{ID}} & \textcolor{white}{\textbf{C1}} & \textcolor{white}{\textbf{C2}} & \textcolor{white}{\textbf{C3}} & \textcolor{white}{\textbf{C4}} & \textcolor{white}{\textbf{C5}} & \textcolor{white}{\textbf{Estado}} \\ \hline
    \endfirsthead
    
    % --- Esta cabecera aparecerá automáticamente si la tabla salta de página ---
    \multicolumn{7}{c}{{\bfseries \dots continuación de la página anterior}} \\
    \hline
    \rowcolor{UCO}
    \textcolor{white}{\textbf{ID}} & \textcolor{white}{\textbf{C1}} & \textcolor{white}{\textbf{C2}} & \textcolor{white}{\textbf{C3}} & \textcolor{white}{\textbf{C4}} & \textcolor{white}{\textbf{C5}} & \textcolor{white}{\textbf{Estado}} \\ \hline
    \endhead

    % --- CONTENIDO DE LA TABLA ---
    \multicolumn{7}{|c|}{\cellcolor{gray!20}\textbf{Módulo 1: Gestión de Usuarios}} \\ \hline
    RF1.1 & \pass & \pass & \pass & \pass & \pass & \textcolor{codegreen}{Aprobado} \\ \hline
    RF1.2 & \pass & \pass & \pass & \pass & \pass & \textcolor{codegreen}{Aprobado} \\ \hline
    RF1.3 & \pass & \pass & \pass & \pass & \pass & \textcolor{codegreen}{Aprobado} \\ \hline
    RF1.4 & \pass & \pass & \pass & \pass & \pass & \textcolor{codegreen}{Aprobado} \\ \hline
    RF1.5 & \pass & \pass & \pass & \pass & \pass & \textcolor{codegreen}{Aprobado} \\ \hline

    \multicolumn{7}{|c|}{\cellcolor{gray!20}\textbf{Módulo 2: Registro Diario y Etiquetas}} \\ \hline
    RF2.1 & \pass & \pass & \pass & \pass & \pass & \textcolor{codegreen}{Aprobado} \\ \hline
    RF2.2 & \pass & \pass & \pass & \pass & \fail & \textcolor{orange}{A refinar} \\ \hline
    RF2.3 & \pass & \pass & \pass & \pass & \pass & \textcolor{codegreen}{Aprobado} \\ \hline
    RF2.4 & \pass & \pass & \pass & \pass & \pass & \textcolor{codegreen}{Aprobado} \\ \hline
    RF2.5 & \pass & \pass & \pass & \pass & \pass & \textcolor{codegreen}{Aprobado} \\ \hline
    RF2.6 & \pass & \pass & \pass & \pass & \pass & \textcolor{codegreen}{Aprobado} \\ \hline
    RF2.7 & \pass & \pass & \pass & \pass & \pass & \textcolor{codegreen}{Aprobado} \\ \hline
    RF2.8 & \pass & \pass & \pass & \pass & \pass & \textcolor{codegreen}{Aprobado} \\ \hline
    RF2.9 & \fail & \pass & \pass & \pass & \fail & \textcolor{orange}{A refinar} \\ \hline
    RF2.10 & \pass & \pass & \pass & \pass & \pass & \textcolor{codegreen}{Aprobado} \\ \hline

    \multicolumn{7}{|c|}{\cellcolor{gray!20}\textbf{Módulo 3: Panel Principal e Historial}} \\ \hline
    RF3.1 & \pass & \pass & \pass & \pass & \pass & \textcolor{codegreen}{Aprobado} \\ \hline
    RF3.2 & \fail & \fail & \pass & \pass & \pass & \textcolor{orange}{A refinar} \\ \hline
    RF3.3 & \pass & \pass & \pass & \pass & \pass & \textcolor{codegreen}{Aprobado} \\ \hline
    RF3.4 & \pass & \pass & \pass & \pass & \pass & \textcolor{codegreen}{Aprobado} \\ \hline
    RF3.5 & \pass & \pass & \pass & \pass & \pass & \textcolor{codegreen}{Aprobado} \\ \hline
    RF3.6 & \pass & \pass & \pass & \pass & \pass & \textcolor{codegreen}{Aprobado} \\ \hline

    \multicolumn{7}{|c|}{\cellcolor{gray!20}\textbf{Módulo 4: Análisis de Datos (Estadísticas)}} \\ \hline
    RF4.1 & \pass & \pass & \pass & \pass & \pass & \textcolor{codegreen}{Aprobado} \\ \hline
    RF4.2 & \pass & \pass & \pass & \pass & \pass & \textcolor{codegreen}{Aprobado} \\ \hline
    RF4.3 & \fail & \fail & \pass & \pass & \pass & \textcolor{orange}{A refinar} \\ \hline
    RF4.4 & \pass & \pass & \pass & \pass & \pass & \textcolor{codegreen}{Aprobado} \\ \hline
    RF4.5 & \pass & \pass & \pass & \pass & \pass & \textcolor{codegreen}{Aprobado} \\ \hline
    RF4.6 & \pass & \pass & \pass & \pass & \pass & \textcolor{codegreen}{Aprobado} \\ \hline
    RF4.7 & \pass & \pass & \pass & \pass & \pass & \textcolor{codegreen}{Aprobado} \\ \hline
    RF4.8 & \pass & \pass & \pass & \pass & \pass & \textcolor{codegreen}{Aprobado} \\ \hline

    \multicolumn{7}{|c|}{\cellcolor{gray!20}\textbf{Módulo 5: Personalización y Sistema}} \\ \hline
    RF5.1 & \pass & \pass & \pass & \pass & \pass & \textcolor{codegreen}{Aprobado} \\ \hline
    RF5.2 & \pass & \pass & \pass & \pass & \pass & \textcolor{codegreen}{Aprobado} \\ \hline
    RF5.3 & \fail & \pass & \pass & \pass & \pass & \textcolor{orange}{A refinar} \\ \hline
    RF5.4 & \pass & \pass & \pass & \pass & \pass & \textcolor{codegreen}{Aprobado} \\ \hline
    RF5.5 & \pass & \pass & \pass & \pass & \pass & \textcolor{codegreen}{Aprobado} \\ \hline
    \caption{Matriz de Verificación de Requisitos Funcionales} \label{tab:matriz_rf} \\

\end{xltabular}

\subsubsection{Matriz de Verificación: Requisitos No Funcionales}

\renewcommand{\arraystretch}{1.3}
\begin{xltabular}{\textwidth}{|c|>{\centering\arraybackslash}X|>{\centering\arraybackslash}X|>{\centering\arraybackslash}X|>{\centering\arraybackslash}X|>{\centering\arraybackslash}X|c|}
    \hline
    \rowcolor{UCO}
    \textcolor{white}{\textbf{ID}} & \textcolor{white}{\textbf{C1}} & \textcolor{white}{\textbf{C2}} & \textcolor{white}{\textbf{C3}} & \textcolor{white}{\textbf{C4}} & \textcolor{white}{\textbf{C5}} & \textcolor{white}{\textbf{Estado}} \\ \hline
    \endfirsthead
    
    % --- Cabecera para saltos de página ---
    \multicolumn{7}{c}{{\bfseries \dots continuación de la página anterior}} \\
    \hline
    \rowcolor{UCO}
    \textcolor{white}{\textbf{ID}} & \textcolor{white}{\textbf{C1}} & \textcolor{white}{\textbf{C2}} & \textcolor{white}{\textbf{C3}} & \textcolor{white}{\textbf{C4}} & \textcolor{white}{\textbf{C5}} & \textcolor{white}{\textbf{Estado}} \\ \hline
    \endhead


    % --- CONTENIDO DE LA TABLA ---
    \multicolumn{7}{|c|}{\cellcolor{gray!20}\textbf{Bloque 1: Compatibilidad y Portabilidad}} \\ \hline
    RNF1.1 & \pass & \pass & \pass & \pass & \pass & \textcolor{codegreen}{Aprobado} \\ \hline
    RNF1.2 & \pass & \pass & \pass & \pass & \pass & \textcolor{codegreen}{Aprobado} \\ \hline
    RNF1.3 & \pass & \pass & \pass & \pass & \pass & \textcolor{codegreen}{Aprobado} \\ \hline

    \multicolumn{7}{|c|}{\cellcolor{gray!20}\textbf{Bloque 2: Rendimiento y Disponibilidad}} \\ \hline
    RNF2.1 & \pass & \pass & \pass & \pass & \pass & \textcolor{codegreen}{Aprobado} \\ \hline
    RNF2.2 & \pass & \pass & \pass & \pass & \pass & \textcolor{codegreen}{Aprobado} \\ \hline
    RNF2.3 & \pass & \pass & \pass & \pass & \pass & \textcolor{codegreen}{Aprobado} \\ \hline
    RNF2.4 & \pass & \pass & \pass & \pass & \pass & \textcolor{codegreen}{Aprobado} \\ \hline

    \multicolumn{7}{|c|}{\cellcolor{gray!20}\textbf{Bloque 3: Usabilidad e Interfaz}} \\ \hline
    RNF3.1 & \pass & \pass & \pass & \pass & \pass & \textcolor{codegreen}{Aprobado} \\ \hline
    RNF3.2 & \pass & \pass & \pass & \pass & \pass & \textcolor{codegreen}{Aprobado} \\ \hline
    RNF3.3 & \fail & \fail & \pass & \pass & \pass & \textcolor{orange}{A refinar} \\ \hline

    \multicolumn{7}{|c|}{\cellcolor{gray!20}\textbf{Bloque 4: Seguridad y Privacidad}} \\ \hline
    RNF4.1 & \pass & \pass & \pass & \pass & \pass & \textcolor{codegreen}{Aprobado} \\ \hline
    RNF4.2 & \pass & \pass & \pass & \pass & \pass & \textcolor{codegreen}{Aprobado} \\ \hline
    RNF4.3 & \pass & \pass & \pass & \pass & \pass & \textcolor{codegreen}{Aprobado} \\ \hline
    RNF4.4 & \pass & \pass & \pass & \pass & \pass & \textcolor{codegreen}{Aprobado} \\ \hline

    \caption{Matriz de Verificación de Requisitos No Funcionales} \label{tab:matriz_rnf} \\

\end{xltabular}

\subsubsection{Matriz de Verificación: Requisitos de Información}

\renewcommand{\arraystretch}{1.3}
\begin{xltabular}{\textwidth}{|c|>{\centering\arraybackslash}X|>{\centering\arraybackslash}X|>{\centering\arraybackslash}X|>{\centering\arraybackslash}X|>{\centering\arraybackslash}X|c|}
    \hline
    \rowcolor{UCO}
    \textcolor{white}{\textbf{ID}} & \textcolor{white}{\textbf{C1}} & \textcolor{white}{\textbf{C2}} & \textcolor{white}{\textbf{C3}} & \textcolor{white}{\textbf{C4}} & \textcolor{white}{\textbf{C5}} & \textcolor{white}{\textbf{Estado}} \\ \hline
    \endfirsthead
    
    % --- Cabecera para saltos de página ---
    \multicolumn{7}{c}{{\bfseries \dots continuación de la página anterior}} \\
    \hline
    \rowcolor{UCO}
    \textcolor{white}{\textbf{ID}} & \textcolor{white}{\textbf{C1}} & \textcolor{white}{\textbf{C2}} & \textcolor{white}{\textbf{C3}} & \textcolor{white}{\textbf{C4}} & \textcolor{white}{\textbf{C5}} & \textcolor{white}{\textbf{Estado}} \\ \hline
    \endhead


    % --- CONTENIDO DE LA TABLA ---
    RI-01 (Usuario) & \pass & \pass & \pass & \pass & \pass & \textcolor{codegreen}{Aprobado} \\ \hline
    RI-02 (Etiqueta) & \pass & \pass & \pass & \pass & \pass & \textcolor{codegreen}{Aprobado} \\ \hline
    RI-03 (Registro) & \pass & \pass & \pass & \pass & \pass & \textcolor{codegreen}{Aprobado} \\ \hline
    \caption{Matriz de Verificación de Requisitos de Información} \label{tab:matriz_ri} \\

\end{xltabular}

\subsubsection{Refinamiento y Corrección de Requisitos}

Tras la aplicación de la matriz de verificación de requisitos se ha detectado que seis requisitos(marcados en la matriz con estado \textit{''A refinar''}) no han superado los criterios de calidad establecidos. Por tanto, para garantizar la robustez en la definición de requisitos, estos han sido sometidos a un proceso de reescritura, para asegurar que cumplen correctamente con los criterios de calidad: 


\begin{itemize}
    \item \textbf{RF2.2 (Fallo en C5 - Incompleto):}
    \begin{itemize}
        \item \textit{Original:} El sistema establecerá la fecha actual por defecto al iniciar la creación de un nuevo registro, permitiendo su modificación a fechas pasadas.
        \item \textit{Corregido:} El sistema establecerá la fecha actual por defecto al iniciar la creación de un nuevo registro, permitiendo su modificación únicamente a fechas pasadas y bloqueando explícitamente la selección de fechas futuras.
    \end{itemize}

    \item \textbf{RF2.9 (Fallo en C1 y C5 - Subjetivo e incompleto):}
    \begin{itemize}
        \item \textit{Original:} El sistema permitirá al usuario modificar cualquier dato de un registro propio previamente guardado.
        \item \textit{Corregido:} El sistema permitirá al usuario editar los atributos modificables (puntuación global, fecha asignada, comentario y etiquetas vinculadas) de un registro propio, aplicando las mismas reglas de validación y obligatoriedad definidas para la creación del registro.
    \end{itemize}
    
    \item \textbf{RF3.2 (Fallo en C1 y C2 - Subjetivo e inmedible):}
    \begin{itemize}
        \item \textit{Original:} El sistema proporcionará en el panel principal un acceso directo y visualmente destacado para la creación de un nuevo registro diario.
        \item \textit{Corregido:} El sistema proporcionará en el panel principal un botón de acceso directo a la creación de un nuevo registro, accesible sin necesidad de navegación adicional.
    \end{itemize}

    \item \textbf{RF4.3 (Fallo en C1 - Ambiguo):}
    \begin{itemize}
        \item \textit{Original:} El sistema calculará y mostrará qué etiquetas presentan el mayor impacto positivo y negativo sobre el estado de ánimo global del usuario.
        \item \textit{Corregido:} El sistema calculará y mostrará qué etiquetas están asociadas matemáticamente con las puntuaciones globales medias más altas (impacto positivo) y más bajas (impacto negativo) del estado de ánimo del usuario.
    \end{itemize}

    \item \textbf{RF5.3 (Fallo en C1 - Ambiguo):}
    \begin{itemize}
        \item \textit{Original:} El sistema mostrará en el panel principal un mensaje de bienvenida dinámico, adaptado contextualmente a factores como la hora del día o la actividad reciente.
        \item \textit{Corregido:} El sistema mostrará en el panel principal un mensaje de bienvenida dinámico determinado por dos parámetros: la franja horaria local (mañana, tarde o noche) y el estado de su registro diario (pendiente de rellenar o ya completado en el día actual).
    \end{itemize}

    \item \textbf{RNF3.3 (Fallo en C1 y C2 - Subjetivo e inmedible):}
    \begin{itemize}
        \item \textit{Original:} El sistema proporcionará retroalimentación visual inmediata ante cualquier acción importante (ej. notificaciones de éxito o alertas claras de error).
        \item \textit{Corregido:} El sistema proporcionará retroalimentación visual en un tiempo máximo de 0,5 segundos ante cualquier acción importante (ej. notificaciones de éxito o alertas claras de error).
    \end{itemize}
\end{itemize}

\textit{Nota: Para mantener la coherencia del documento, las tablas generales de especificación de requisitos presentadas en los apartados anteriores ya reflejan estas versiones finales y corregidas.}

\newpage

\subsection{Estudio de Viabilidad Técnica y Selección Tecnológica}

Para asegurar que MoodNest cumpla con los requisitos de rendimiento, seguridad y escalabilidad definidos, se ha realizado un estudio de las alternativas tecnológicas disponibles. La elección de la arquitectura no responde a una preferencia personal, sino a la necesidad de cubrir las demandas específicas de la gestión de datos y el análisis estadístico del sistema.

\subsubsection{Modelo de Datos}
El núcleo de MoodNest es el ''Registro Diario''. Al analizar los requisitos del Módulo 2, se observa que la información de cada registro no es rígida: un usuario puede asociar dos etiquetas a un día y diez etiquetas al siguiente, añadir comentarios de longitud variable o incluir puntuaciones específicas solo para ciertos factores.

En un modelo relacional (SQL) como PostgreSQL o MySQL, esta flexibilidad obligaría a realizar diseños de tablas complejos con múltiples uniones o el uso extensivo de valores nulos, lo que penalizaría el rendimiento de las consultas (incumpliendo el RNF2.1). Por ello, se ha optado por un modelo NoSQL orientado a documentos con MongoDB. Esta tecnología permite almacenar cada registro como un documento JSON independiente, facilitando que la estructura de los datos evolucione de forma natural y que el acceso a la información sea directo y mucho más veloz.

\subsubsection{Lógica de Negocio y Seguridad}
Para la construcción del servidor, se evaluaron alternativas como Node.js y Spring Boot. Aunque Node.js ofrece una gran rapidez de ejecución, la naturaleza de MoodNest (que maneja datos personales sensibles) requiere un entorno con un enfoque robusto en seguridad y estabilidad.

Se ha seleccionado Spring Boot (Java) debido a su madurez y a la potencia de su ecosistema. Especialmente crítico es el uso de Spring Security, que permite implementar de forma nativa la protección de datos y el acceso mediante tokens JWT exigidos en el RNF4.2. Además, la integración con MongoDB a través de Spring Data permite que el desarrollo sea mucho más ágil, garantizando que el backend sea capaz de procesar las peticiones concurrentes definidas en el RNF2.3 sin degradar la experiencia del usuario.

\subsubsection{Interfaz de Usuario y Visualización de Datos}
La interfaz de la aplicación tiene dos retos principales: la agilidad (RNF3.2) y la complejidad de los gráficos estadísticos (Módulo 4). Tras analizar opciones como Angular o Vue, se ha determinado que React es la opción más viable para este proyecto, utilizando la biblioteca Chart.js para la generación y gestión de todas las visualizaciones del sistema.

React permite construir la interfaz mediante componentes reutilizables (calendarios, selectores de humor, widgets de estadísticas), lo que facilita enormemente el mantenimiento del código. Su capacidad para gestionar el estado de la aplicación de forma reactiva asegura que los cambios en los datos se reflejen instantáneamente en los gráficos generados con Chart.js sin necesidad de recargar la página, cumpliendo así con los tiempos de renderizado exigidos. Asimismo, la integración de esta biblioteca específica garantiza que la implementación de los gráficos de líneas, barras y circulares sea técnica y operativamente posible.

Como conclusión de este análisis, se determina que el stack tecnológico seleccionado es capaz de soportar todas las funcionalidades descritas y garantiza el cumplimiento de los requisitos no funcionales de seguridad y rendimiento. Por tanto, el desarrollo de MoodNest se considera técnicamente viable bajo esta arquitectura, permitiendo avanzar con seguridad hacia la fase de especificación funcional del sistema.



\newpage
\subsection{Análisis Funcional}

En esta sección se traduce el catálogo de requisitos en interacciones narrativas. Se define qué actores estarán presentes en el sistema y qué acciones pueden realizar, asegurando que toda la funcionalidad técnica tenga una representación orientada al usuario.

\subsubsection{Identificación y Descripción de Actores}

Un actor representa un rol que un usuario desempeña al interactuar con el sistema. En MoodNest se han identificado dos actores:


\begin{table}[h!]
    \centering
    \renewcommand{\arraystretch}{1.5}
    \begin{tabularx}{\textwidth}{|c|X|}
        \hline
        \rowcolor{UCO} 
        \textcolor{white}{\textbf{Actor}} & \multicolumn{1}{c|}{\textcolor{white}{\textbf{Descripción}}} \\ \hline
        Usuario Registrado & Actor principal del sistema que ha superado el proceso de autenticación y tiene acceso total a la gestión de sus registros, etiquetas y análisis estadístico. \\ \hline
        Usuario Visitante & Usuario no autenticado que accede a la plataforma por primera vez o no ha iniciado sesión. Limitado a las funciones de registro e inicio de sesión. \\ \hline
    \end{tabularx}
    \caption{Descripción de los actores}
    \label{tab:actores}
\end{table}

\subsubsection{Definición de Casos de Uso}

A continuación, se presenta el catálogo general de Casos de Uso del sistema, organizados por su finalidad funcional.

% Usamos >{\raggedright\arraybackslash} para que no intente justificar y no de errores de hboxes
\renewcommand{\arraystretch}{1.4}
\begin{xltabular}{\textwidth}{|c|>{\raggedright\arraybackslash}p{3.5cm}|>{\raggedright\arraybackslash}X|c|}    \hline
    \rowcolor{UCO} 
    \textcolor{white}{\textbf{Id}} & \multicolumn{1}{c|}{\textcolor{white}{\textbf{Nombre}}} & \multicolumn{1}{c|}{\textcolor{white}{\textbf{Descripción}}} & \textcolor{white}{\textbf{Actor}} \\ \hline
    \endfirsthead
    
    \multicolumn{4}{c}{{\bfseries \dots continuación de la página anterior}} \\
    \hline
    \rowcolor{UCO} 
    \textcolor{white}{\textbf{Id}} & \multicolumn{1}{c|}{\textcolor{white}{\textbf{Nombre}}} & \multicolumn{1}{c|}{\textcolor{white}{\textbf{Descripción}}} & \textcolor{white}{\textbf{Actor}} \\ \hline
    \endhead

    % --- Grupo 1 ---
    \multicolumn{4}{|c|}{\cellcolor{gray!20}\textbf{Gestión de Cuenta}} \\ \hline
    CU1 & Registrar Usuario & Un usuario se registra en MoodNest proporcionando nombre, correo y contraseña. & Visitante \\ \hline
    CU2 & Iniciar Sesión & Un usuario accede al sistema con su correo y contraseña. & Visitante \\ \hline
    CU3 & Gestionar Perfil & El usuario modifica su nombre de visualización o su contraseña. & Registrado \\ \hline
    CU4 & Cerrar Sesión & El usuario finaliza su sesión de forma segura. & Registrado \\ \hline
    CU5 & Eliminar Cuenta & El usuario borra permanentemente su cuenta y todos sus datos. & Registrado \\ \hline
    
    % --- Grupo 2 ---
    \multicolumn{4}{|c|}{\cellcolor{gray!20}\textbf{Gestión de Registros y Etiquetas}} \\ \hline
    CU6 & Crear Registro Diario & El usuario crea una nueva entrada diaria con sus etiquetas y puntuaciones. & Registrado \\ \hline
    CU7 & Gestionar Registro & El usuario edita o elimina uno de sus registros diarios. & Registrado \\ \hline
    CU8 & Gestionar Etiqueta & El usuario crea, edita o elimina una etiqueta de su catálogo personal. & Registrado \\ \hline
    
    % --- Grupo 3 ---
    \multicolumn{4}{|c|}{\cellcolor{gray!20}\textbf{Visualización y Consulta}} \\ \hline
    CU9 & Visualizar Dashboard & El usuario consulta el panel principal con el resumen de su actividad. & Registrado \\ \hline
    CU10 & Visualizar Historial & El usuario consulta sus registros pasados a través de un calendario. & Registrado \\ \hline
    
    % --- Grupo 4 ---
    \multicolumn{4}{|c|}{\cellcolor{gray!20}\textbf{Análisis Estadístico}} \\ \hline
    CU11 & Consultar Evolución Temporal & El usuario solicita y filtra la gráfica de su estado de ánimo por periodos (semanal, mensual o anual). & Registrado \\ \hline
    CU12 & Visualizar Resumen Estadístico & El usuario accede a la página de análisis y visualiza distintas estadísticas sobre sus registros diarios y etiquetas. & Registrado \\ \hline
    CU13 & Analizar Impacto de Etiquetas & El usuario selecciona etiquetas específicas para comparar su influencia en el estado de ánimo global. & Registrado \\ \hline    
    % --- Grupo 5 ---
    \multicolumn{4}{|c|}{\cellcolor{gray!20}\textbf{Personalización del Sistema}} \\ \hline
    CU14 & Personalizar Escala & El usuario configura los textos e iconos de su escala del 1 al 10. & Registrado \\ \hline
    CU15 & Personalizar Interfaz & El usuario selecciona el tema (claro/oscuro) y el color de la aplicación. & Registrado \\ \hline

    \caption{Definición de los casos de uso organizados por módulos}
    \label{tab:definicion_cu}

\end{xltabular}


\subsubsection{Matriz de Trazabilidad}

Con el objetivo de garantizar la integridad y consistencia del diseño funcional, se presenta a continuación la Matriz de Trazabilidad, la cual nos permite realizar una validación cruzada entre los Requisitos Funcionales y Casos de Uso.

La elaboración de esta matriz tiene como objetivo confirmar que cada uno de los requisitos detectados está satisfecho por, al menos, un caso de uso, evitando que queden necesidades del usuario sin una solución técnica y a su vez garantizar que cada caso de uso existe para dar cumplimiento a un requisito específico, evitando la implementación de funcionalidades innecesarias que no aporten valor real al proyecto.

% Definimos un tipo de columna X centrada para que la matriz quede simétrica
\newcolumntype{Y}{>{\centering\arraybackslash}X}

\begin{center}
\footnotesize 
\renewcommand{\arraystretch}{1.2}
% Usamos Y para las 15 columnas de los CUs para que se estiren hasta ocupar \textwidth
\begin{xltabular}{\textwidth}{|c|Y|Y|Y|Y|Y|Y|Y|Y|Y|Y|Y|Y|Y|Y|Y|}
    \hline
    \rowcolor{UCO}
    \textcolor{white}{\textbf{ID RF}} & 
    \rotatebox{90}{\textcolor{white}{CU1}} & \rotatebox{90}{\textcolor{white}{CU2}} & \rotatebox{90}{\textcolor{white}{CU3}} & \rotatebox{90}{\textcolor{white}{CU4}} & \rotatebox{90}{\textcolor{white}{CU5}} & 
    \rotatebox{90}{\textcolor{white}{CU6}} & \rotatebox{90}{\textcolor{white}{CU7}} & \rotatebox{90}{\textcolor{white}{CU8}} & 
    \rotatebox{90}{\textcolor{white}{CU9}} & \rotatebox{90}{\textcolor{white}{CU10}} & 
    \rotatebox{90}{\textcolor{white}{CU11}} & \rotatebox{90}{\textcolor{white}{CU12}} & \rotatebox{90}{\textcolor{white}{CU13}} & 
    \rotatebox{90}{\textcolor{white}{CU14}} & \rotatebox{90}{\textcolor{white}{CU15}} \\ \hline
    \endfirsthead

    \hline
    \rowcolor{UCO}
    \textcolor{white}{\textbf{ID RF}} & 
    \rotatebox{90}{\textcolor{white}{CU1}} & \rotatebox{90}{\textcolor{white}{CU2}} & \rotatebox{90}{\textcolor{white}{CU3}} & \rotatebox{90}{\textcolor{white}{CU4}} & \rotatebox{90}{\textcolor{white}{CU5}} & 
    \rotatebox{90}{\textcolor{white}{CU6}} & \rotatebox{90}{\textcolor{white}{CU7}} & \rotatebox{90}{\textcolor{white}{CU8}} & 
    \rotatebox{90}{\textcolor{white}{CU9}} & \rotatebox{90}{\textcolor{white}{CU10}} & 
    \rotatebox{90}{\textcolor{white}{CU11}} & \rotatebox{90}{\textcolor{white}{CU12}} & \rotatebox{90}{\textcolor{white}{CU13}} & 
    \rotatebox{90}{\textcolor{white}{CU14}} & \rotatebox{90}{\textcolor{white}{CU15}} \\ \hline
    \endhead


    \multicolumn{16}{|c|}{\cellcolor{gray!20}\textbf{Módulo 1: Gestión de Usuarios}} \\ \hline
    RF1.1 & X &   &   &   &   &   &   &   &   &   &   &   &   &   &   \\ \hline
    RF1.2 &   & X &   &   &   &   &   &   &   &   &   &   &   &   &   \\ \hline
    RF1.3 &   &   & X &   &   &   &   &   &   &   &   &   &   &   &   \\ \hline
    RF1.4 &   &   &   & X &   &   &   &   &   &   &   &   &   &   &   \\ \hline
    RF1.5 &   &   &   &   & X &   &   &   &   &   &   &   &   &   &   \\ \hline

    \multicolumn{16}{|c|}{\cellcolor{gray!20}\textbf{Módulo 2: Registro Diario y Etiquetas}} \\ \hline
    RF2.1 &   &   &   &   &   & X &   &   &   &   &   &   &   &   &   \\ \hline
    RF2.2 &   &   &   &   &   & X &   &   &   &   &   &   &   &   &   \\ \hline
    RF2.3 &   &   &   &   &   & X &   &   &   &   &   &   &   &   &   \\ \hline
    RF2.4 &   &   &   &   &   & X &   &   &   &   &   &   &   &   &   \\ \hline
    RF2.5 &   &   &   &   &   &   &   & X &   &   &   &   &   &   &   \\ \hline
    RF2.6 &   &   &   &   &   &   &   & X &   &   &   &   &   &   &   \\ \hline
    RF2.7 &   &   &   &   &   & X &   &   &   &   &   &   &   &   &   \\ \hline
    RF2.8 &   &   &   &   &   & X &   &   &   &   &   &   &   &   &   \\ \hline
    RF2.9 &   &   &   &   &   &   & X &   &   &   &   &   &   &   &   \\ \hline
    RF2.10&   &   &   &   &   &   & X &   &   &   &   &   &   &   &   \\ \hline

    \multicolumn{16}{|c|}{\cellcolor{gray!20}\textbf{Módulo 3: Panel Principal e Historial}} \\ \hline
    RF3.1 &   &   &   &   &   &   &   &   & X &   &   &   &   &   &   \\ \hline
    RF3.2 &   &   &   &   &   &   &   &   & X &   &   &   &   &   &   \\ \hline
    RF3.3 &   &   &   &   &   &   &   &   & X &   &   &   &   &   &   \\ \hline
    RF3.4 &   &   &   &   &   &   &   &   & X &   &   &   &   &   &   \\ \hline
    RF3.5 &   &   &   &   &   &   &   &   &   & X &   &   &   &   &   \\ \hline
    RF3.6 &   &   &   &   &   &   &   &   &   & X &   &   &   &   &   \\ \hline

    \multicolumn{16}{|c|}{\cellcolor{gray!20}\textbf{Módulo 4: Análisis de Datos (Estadísticas)}} \\ \hline
    RF4.1 &   &   &   &   &   &   &   &   &   &   & X &   &   &   &   \\ \hline
    RF4.2 &   &   &   &   &   &   &   &   &   &   & X &   &   &   &   \\ \hline
    RF4.3 &   &   &   &   &   &   &   &   &   &   &   & X &   &   &   \\ \hline
    RF4.4 &   &   &   &   &   &   &   &   &   &   &   &   & X &   &   \\ \hline
    RF4.5 &   &   &   &   &   &   &   &   &   &   &   &   & X &   &   \\ \hline
    RF4.6 &   &   &   &   &   &   &   &   &   &   &   &   & X &   &   \\ \hline
    RF4.7 &   &   &   &   &   &   &   &   &   &   &   & X &   &   &   \\ \hline
    RF4.8 &   &   &   &   &   &   &   &   &   &   &   & X &   &   &   \\ \hline

    \multicolumn{16}{|c|}{\cellcolor{gray!20}\textbf{Módulo 5: Personalización y Sistema}} \\ \hline
    RF5.1 &   &   &   &   &   &   &   &   &   &   &   &   &   & X &   \\ \hline
    RF5.2 &   &   &   &   &   &   &   &   &   &   &   &   &   & X &   \\ \hline
    RF5.3 &   &   &   &   &   &   &   &   & X &   &   &   &   &   &   \\ \hline
    RF5.4 &   &   &   &   &   &   &   &   &   &   &   &   &   &   & X \\ \hline
    RF5.5 &   &   &   &   &   &   &   &   &   &   &   &   &   &   & X \\ \hline


    \caption{Matriz de validación cruzada de Requisitos y Casos de Uso} \label{tab:matriz_trazabilidad} \\
\end{xltabular}
\end{center}

\newpage
\subsubsection{Diagrama General de Casos de Uso}

A continuación, en la Figura \ref{fig:diagrama_cu}, se presenta el diagrama general de Casos de Uso de la aplicación. Este modelo visualiza la frontera del sistema y establece las relaciones estructurales entre los actores identificados y las funcionalidades principales, incluyendo las dependencias de inclusión (\textit{include}) y extensión (\textit{extend}) que rigen la lógica de navegación.

\begin{figure}[H]
    \centering
    % Ajusta el width si ves que la imagen queda muy grande (ej. width=0.9\textwidth)
    \includegraphics[width=\textwidth]{diagrama_CU.pdf}
    \caption{Diagrama General de Casos de Uso de MoodNest}
    \label{fig:diagrama_cu}
\end{figure}


\newpage
\subsubsection{Especificación Detallada de Casos de Uso}

\begingroup
\small 
\renewcommand{\arraystretch}{1.4}
% Añadido \columncolor{gray!20} a la primera columna
\begin{xltabular}{\textwidth}{|>{\columncolor{gray!20}\bfseries}p{3.5cm}|X|}
    
    \rowcolor{UCO}
    \textcolor{white}{Campo} & \textcolor{white}{Descripción} \\ \hline
    \endfirsthead

    \multicolumn{2}{c}{{\bfseries \dots continuación de la página anterior}} \\
    \hline
    \rowcolor{UCO}
    \textcolor{white}{Campo} & \textcolor{white}{Descripción} \\ \hline
    \endhead


    % CONTENIDO DE LA TABLA
    Nombre & Registrar Usuario \\ \hline
    Identificador & CU1 \\ \hline
    Descripción & Un usuario se registra en MoodNest proporcionando su nombre, correo electrónico y contraseña. \\ \hline
    Actor Principal & Usuario Visitante \\ \hline
    Actor Secundario & Ninguno \\ \hline
    Precondición & Ninguna \\ \hline
    Flujo Principal & 
    \begin{enumerate}[leftmargin=0.5cm, topsep=0pt, partopsep=0pt, parsep=0pt, itemsep=4pt]
        \item El caso de uso comienza cuando un usuario selecciona la opción ''Registrarse''.
        \item El sistema solicita las credenciales y datos necesarios (nombre, correo electrónico y contraseña).
        \item El usuario introduce la información requerida.
        \item El sistema comprueba que la información introducida es válida y que el correo no está registrado previamente en la base de datos.
        \item El sistema cifra la contraseña y registra un nuevo usuario.
        \item El sistema redirige a la pantalla de dashboard del usuario.
    \end{enumerate} 
    \\ \hline
    Flujo Alternativo & 
    \begin{enumerate}[leftmargin=0.5cm, topsep=0pt, partopsep=0pt, parsep=0pt, itemsep=4pt]
        \item \textbf{Datos Inválidos o Duplicados}: Se activa en el paso 4 si la información no cumple con el formato exigido o el correo ya existe en la base de datos. El sistema muestra un mensaje de alerta y solicita que se introduzcan de nuevo.
        \item \textbf{Cancelar}: Comienza en cualquier punto del flujo principal. El usuario selecciona cancelar o vuelve atrás. El sistema cancela el proceso en curso y regresa a la pantalla de inicio.
    \end{enumerate}
    \\ \hline
    Postcondición & El usuario queda registrado en el sistema y listo para iniciar sesión. \\ \hline
    \caption{CU1: Registrar Usuario}
    \label{tab:esp_cu1} \\
\end{xltabular}
\endgroup


\begingroup
\small 
\renewcommand{\arraystretch}{1.4}
\begin{xltabular}{\textwidth}{|>{\columncolor{gray!20}\bfseries}p{3.5cm}|X|}
    \rowcolor{UCO}
    \textcolor{white}{Campo} & \textcolor{white}{Descripción} \\ \hline
    \endfirsthead

    \multicolumn{2}{c}{{\bfseries \dots continuación de la página anterior}} \\
    \hline
    \rowcolor{UCO}
    \textcolor{white}{Campo} & \textcolor{white}{Descripción} \\ \hline
    \endhead

    % CONTENIDO DE LA TABLA
    Nombre & Iniciar Sesión \\ \hline
    Identificador & CU2 \\ \hline
    Descripción & Un usuario accede al sistema introduciendo sus credenciales para obtener acceso a sus datos privados. \\ \hline
    Actor Principal & Usuario Visitante \\ \hline
    Actor Secundario & Ninguno \\ \hline
    Precondición & El usuario debe tener una cuenta previamente registrada en el sistema. \\ \hline
    Flujo Principal & 
    \begin{enumerate}[leftmargin=0.5cm, topsep=0pt, partopsep=0pt, parsep=0pt, itemsep=4pt]
        \item El caso de uso comienza cuando un usuario accede a la pantalla de inicio de sesión.
        \item El sistema solicita el correo electrónico y la contraseña.
        \item El usuario introduce sus credenciales y selecciona ''Iniciar Sesión''.
        \item El sistema verifica que la cuenta existe y que la contraseña es correcta.
        \item El sistema genera un token de autenticación válido para la sesión.
        \item El sistema redirige al usuario al Panel Principal y le concede permisos de ''Usuario Registrado''.
    \end{enumerate} 
    \\ \hline
    Flujo Alternativo & 
    \begin{enumerate}[leftmargin=0.5cm, topsep=0pt, partopsep=0pt, parsep=0pt, itemsep=4pt]
        \item \textbf{Credenciales Incorrectas}: Se activa en el paso 4 si el correo no existe o la contraseña no coincide. El sistema deniega el acceso, muestra un mensaje de aviso y borra el campo de la contraseña por seguridad.
    \end{enumerate}
    \\ \hline
    Postcondición & El usuario pasa a tener el rol de ''Usuario Registrado'' y puede acceder a las funciones principales de la aplicación. \\ \hline
    \caption{CU2: Iniciar Sesión}
    \label{tab:esp_cu2} \\
\end{xltabular}
\endgroup

\begingroup
\small 
\renewcommand{\arraystretch}{1.4}
\begin{xltabular}{\textwidth}{|>{\columncolor{gray!20}\bfseries}p{3.5cm}|X|}
    \rowcolor{UCO}
    \textcolor{white}{Campo} & \textcolor{white}{Descripción} \\ \hline
    \endfirsthead

    \multicolumn{2}{c}{{\bfseries \dots continuación de la página anterior}} \\
    \hline
    \rowcolor{UCO}
    \textcolor{white}{Campo} & \textcolor{white}{Descripción} \\ \hline
    \endhead

    % CONTENIDO DE LA TABLA
    Nombre & Gestionar Perfil \\ \hline
    Identificador & CU3 \\ \hline
    Descripción & El usuario autenticado modifica sus datos básicos de perfil, como su nombre de visualización o su contraseña de acceso. \\ \hline
    Actor Principal & Usuario Registrado \\ \hline
    Actor Secundario & Ninguno \\ \hline
    Precondición & El usuario debe estar autenticado en el sistema. \\ \hline
    Flujo Principal & 
    \begin{enumerate}[leftmargin=0.5cm, topsep=0pt, partopsep=0pt, parsep=0pt, itemsep=4pt]
        \item El caso de uso comienza cuando el usuario accede a la sección de ''Configuración de Cuenta''.
        \item El sistema recupera y muestra el nombre de visualización actual del usuario y su email.
        \item El usuario modifica su nombre y/o introduce una nueva contraseña en los campos habilitados.
        \item El usuario selecciona la opción "Guardar Cambios".
        \item El sistema valida que los datos introducidos cumplan con los requisitos de formato y seguridad.
        \item El sistema actualiza la información en la base de datos.
        \item El sistema confirma la operación mediante un mensaje de éxito.
    \end{enumerate} 
    \\ \hline
    Flujo Alternativo & 
    \begin{enumerate}[leftmargin=0.5cm, topsep=0pt, partopsep=0pt, parsep=0pt, itemsep=4pt]
        \item \textbf{Validación Fallida}: Si los datos no cumplen con el formato o la contraseña es demasiado débil, el sistema muestra un mensaje de aviso y permite al usuario corregirlos sin abandonar la pantalla.
    \end{enumerate}
    \\ \hline
    Postcondición & Los datos del perfil se actualizan permanentemente. \\ \hline

    \caption{CU3: Gestionar Perfil}
    \label{tab:esp_cu3} \\
    
\end{xltabular}
\endgroup


\begingroup
\small 
\renewcommand{\arraystretch}{1.4}
\begin{xltabular}{\textwidth}{|>{\columncolor{gray!20}\bfseries}p{3.5cm}|X|}
    \rowcolor{UCO}
    \textcolor{white}{Campo} & \textcolor{white}{Descripción} \\ \hline
    \endfirsthead

    \multicolumn{2}{c}{{\bfseries \dots continuación de la página anterior}} \\
    \hline
    \rowcolor{UCO}
    \textcolor{white}{Campo} & \textcolor{white}{Descripción} \\ \hline
    \endhead


    % CONTENIDO DE LA TABLA
    Nombre & Cerrar Sesión \\ \hline
    Identificador & CU4 \\ \hline
    Descripción & El usuario finaliza su sesión activa en el sistema de forma segura para proteger la privacidad de sus datos y registros. \\ \hline
    Actor Principal & Usuario Registrado \\ \hline
    Actor Secundario & Ninguno \\ \hline
    Precondición & El usuario debe tener una sesión activa en el sistema. \\ \hline
    Flujo Principal & 
    \begin{enumerate}[leftmargin=0.5cm, topsep=0pt, partopsep=0pt, parsep=0pt, itemsep=4pt]
        \item El caso de uso comienza cuando el usuario selecciona la opción ''Cerrar Sesión'' o cuando el usuario elimina su cuenta (CU5).
        \item Si la acción es iniciada manualmente por el usuario
            \begin{enumerate}[label=\theenumi.\arabic*, leftmargin=0.6cm, topsep=2pt, itemsep=2pt]
                \item El sistema muestra un cuadro de diálogo para confirmar la salida.
                \item El usuario confirma su intención de salir de la aplicación.
            \end{enumerate}
        \item El sistema elimina el token de autenticación y limpia el estado de la sesión en el almacenamiento local del navegador.
        \item El sistema redirige al usuario a la pantalla pública de bienvenida.
    \end{enumerate}
    \\ \hline
    Flujo Alternativo & 
    \begin{enumerate}[leftmargin=0.5cm, topsep=0pt, partopsep=0pt, parsep=0pt, itemsep=4pt]
        \item \textbf{Cancelar Cierre}: En el paso 2.1, el usuario selecciona ''Cancelar'' o hace clic fuera del cuadro de diálogo. El sistema oculta el aviso y la sesión activa se mantiene intacta.
        \item \textbf{Sesión Expirada}: Si al intentar cerrar la sesión el token del usuario ya había caducado por inactividad, el sistema ignora la confirmación, limpia los datos residuales del navegador y redirige directamente al inicio de sesión.
    \end{enumerate}
    \\ \hline
    Postcondición & La sesión se da por finalizada. El actor pierde el acceso a las funciones privadas y recobra el rol de ''Usuario Visitante''. \\ \hline

    \caption{CU4: Cerrar Sesión}
    \label{tab:esp_cu4} \\
\end{xltabular}
\endgroup


\begingroup
\small 
\renewcommand{\arraystretch}{1.4}
\begin{xltabular}{\textwidth}{|>{\columncolor{gray!20}\bfseries}p{3.5cm}|X|}
    \rowcolor{UCO}
    \textcolor{white}{Campo} & \textcolor{white}{Descripción} \\ \hline
    \endfirsthead

    \multicolumn{2}{c}{{\bfseries \dots continuación de la página anterior}} \\
    \hline
    \rowcolor{UCO}
    \textcolor{white}{Campo} & \textcolor{white}{Descripción} \\ \hline
    \endhead


    % CONTENIDO DE LA TABLA
    Nombre & Eliminar Cuenta \\ \hline
    Identificador & CU5 \\ \hline
    Descripción & El usuario elimina de forma permanente e irreversible su cuenta y todos los datos asociados a la misma (registros, etiquetas y preferencias). \\ \hline
    Actor Principal & Usuario Registrado \\ \hline
    Actor Secundario & Ninguno \\ \hline
    Precondición & El usuario debe estar autenticado en el sistema. \\ \hline
    Flujo Principal & 
    \begin{enumerate}[leftmargin=0.5cm, topsep=0pt, partopsep=0pt, parsep=0pt, itemsep=4pt]
        \item El caso de uso comienza cuando el usuario accede a la configuración de su perfil y selecciona la opción ''Eliminar Cuenta''.
        \item El sistema muestra una advertencia crítica informando que la acción es completamente irreversible.
        \item El sistema solicita al usuario que introduzca su contraseña actual como medida de seguridad adicional.
        \item El usuario introduce su contraseña y confirma la eliminación.
        \item El sistema valida que la contraseña introducida sea correcta.
        \item El sistema procede a eliminar permanentemente el registro del usuario y realiza un borrado en cascada de todos sus registros diarios y etiquetas personalizadas en la base de datos.
        \item El sistema incluye y ejecuta el ''CU3: Cerrar Sesión''.
        \item El sistema muestra un mensaje de confirmación de baja.
    \end{enumerate} 
    \\ \hline
    Flujo Alternativo & 
    \begin{enumerate}[leftmargin=0.5cm, topsep=0pt, partopsep=0pt, parsep=0pt, itemsep=4pt]
        \item \textbf{Contraseña Incorrecta}: Se activa en el paso 5 si la contraseña no coincide con la del usuario. El sistema detiene el proceso de borrado, muestra una alerta de seguridad y permite reintentar la operación.
        \item \textbf{Cancelar Proceso}: En el paso 2 o 3, el usuario selecciona ''Cancelar'' o retrocede. El sistema aborta la operación y mantiene intactos todos los datos y la sesión del usuario.
    \end{enumerate}
    \\ \hline
    Postcondición & Toda la información del usuario desaparece de la plataforma. El actor pierde su cuenta y recobra el rol de ''Usuario Visitante''. \\ \hline

    \caption{CU5: Eliminar Cuenta}
    \label{tab:esp_cu5} \\
    
\end{xltabular}
\endgroup

\begingroup
\small 
\renewcommand{\arraystretch}{1.4}
\begin{xltabular}{\textwidth}{|>{\columncolor{gray!20}\bfseries}p{3.5cm}|X|}
    \rowcolor{UCO}
    \textcolor{white}{Campo} & \textcolor{white}{Descripción} \\ \hline
    \endfirsthead

    \multicolumn{2}{c}{{\bfseries \dots continuación de la página anterior}} \\
    \hline
    \rowcolor{UCO}
    \textcolor{white}{Campo} & \textcolor{white}{Descripción} \\ \hline
    \endhead

    % CONTENIDO DE LA TABLA
    Nombre & Crear Registro Diario \\ \hline
    Identificador & CU6 \\ \hline
    Descripción & El usuario crea una nueva entrada en su historial asignando una fecha, una puntuación global de su estado de ánimo, etiquetas asociadas y un comentario. \\ \hline
    Actor Principal & Usuario Registrado \\ \hline
    Actor Secundario & Ninguno \\ \hline
    Precondición & El usuario debe estar autenticado en el sistema. \\ \hline
    Flujo Principal & 
    \begin{enumerate}[leftmargin=0.5cm, topsep=0pt, partopsep=0pt, parsep=0pt, itemsep=4pt]
        \item El caso de uso comienza cuando el usuario selecciona la opción ''Registrar mi Día'' desde el dashboard.
        \item El sistema presenta el formulario de creación, estableciendo la fecha actual por defecto.
        \item El usuario confirma la fecha actual o selecciona una fecha pasada en el selector.
        \item El usuario selecciona obligatoriamente una puntuación global (del 1 al 10) para definir su estado de ánimo general del día.
        \item El usuario selecciona, de manera opcional, una o múltiples etiquetas de su catálogo personal.
        \item Para cada etiqueta seleccionada, el usuario puede asignarle (opcionalmente) una puntuación específica del 1 al 10.
        \item El usuario escribe, opcionalmente, un texto descriptivo en el campo de comentario.
        \item El usuario selecciona la opción ''Guardar Registro''.
        \item El sistema valida que la fecha sea válida (no futura) y que exista una puntuación global.
        \item El sistema guarda el registro en la base de datos y recalcula la racha de días consecutivos del usuario si procede.
        \item El sistema redirige al usuario a la vista anterior, mostrando un mensaje de éxito.
    \end{enumerate} 
    \\ \hline
    Flujo Alternativo & 
    \begin{enumerate}[leftmargin=0.5cm, topsep=0pt, partopsep=0pt, parsep=0pt, itemsep=4pt]
        \item \textbf{Validación Fallida}: En el paso 9, si los datos no cumplen con el formato o falta por rellenar algún campo, el sistema muestra un mensaje de aviso y permite al usuario corregirlos sin abandonar la pantalla.
        \item \textbf{Cancelar}: En cualquier momento antes de guardar, el usuario selecciona ''Cancelar'' o retrocede. El sistema descarta toda la información introducida y regresa a la pantalla anterior sin realizar cambios en la base de datos.
    \end{enumerate}
    \\ \hline
    Postcondición & El registro se consolida en la base de datos, queda disponible para su consulta en el historial y computa inmediatamente para las gráficas de análisis estadístico. \\ \hline

    \caption{CU6: Crear Registro Diario}
    \label{tab:esp_cu6} \\
\end{xltabular}
\endgroup


\begingroup
\small 
\renewcommand{\arraystretch}{1.4}
\begin{xltabular}{\textwidth}{|>{\columncolor{gray!20}\bfseries}p{3.5cm}|X|}
    \rowcolor{UCO}
    \textcolor{white}{Campo} & \textcolor{white}{Descripción} \\ \hline
    \endfirsthead

    \multicolumn{2}{c}{{\bfseries \dots continuación de la página anterior}} \\
    \hline
    \rowcolor{UCO}
    \textcolor{white}{Campo} & \textcolor{white}{Descripción} \\ \hline
    \endhead

    % CONTENIDO DE LA TABLA
    Nombre & Gestionar Registro \\ \hline
    Identificador & CU7 \\ \hline
    Descripción & El usuario accede a su historial para editar los atributos de un registro existente o para eliminarlo de forma permanente del sistema. \\ \hline
    Actor Principal & Usuario Registrado \\ \hline
    Actor Secundario & Ninguno \\ \hline
    Precondición & El usuario debe estar autenticado y tener al menos un registro previo en su historial. \\ \hline
    Flujo Principal & 
    \begin{enumerate}[leftmargin=0.5cm, topsep=0pt, partopsep=0pt, parsep=0pt, itemsep=4pt]
        \item El caso de uso comienza cuando el usuario selecciona un registro específico desde la vista de Calendario.
        \item El sistema muestra los detalles completos de la entrada y habilita las opciones de ''Editar'' y ''Eliminar''.
        \item El usuario selecciona la opción ''Editar''.
        \item El sistema presenta el formulario del registro con sus datos actuales pre-cargados (fecha, puntuación, etiquetas vinculadas y comentario).
        \item El usuario modifica los atributos deseados.
        \item El usuario selecciona la opción ''Guardar Cambios''.
        \item El sistema valida las modificaciones aplicando las mismas reglas de negocio y obligatoriedad que en la creación de registros.
        \item El sistema actualiza el registro en la base de datos y recalcula las estadísticas o la racha si fuese necesario.
        \item El sistema redirige a la vista del historial, mostrando un mensaje de éxito.
    \end{enumerate} 
    \\ \hline
    Flujo Alternativo & 
    \begin{enumerate}[leftmargin=0.5cm, topsep=0pt, partopsep=0pt, parsep=0pt, itemsep=4pt]
        \item \textbf{Eliminar Registro}: En el paso 3, el usuario selecciona ''Eliminar''. El sistema muestra un cuadro de diálogo advirtiendo que la acción es irreversible. El usuario confirma. El sistema borra el registro permanentemente, recalcula la racha/estadísticas afectadas y redirige a la pantalla anterior con un mensaje de éxito.
        \item \textbf{Errores de Validación}: Se activa en el paso 7 si los datos modificados no cumplen con el formato. El sistema detiene el proceso y muestra un mensaje de error.
        \item \textbf{Cancelar}: En cualquier momento durante la edición o en el cuadro de confirmación de eliminación, el usuario selecciona ''Cancelar'' o retrocede. El sistema cancela el proceso y mantiene el registro original sin alteraciones.
    \end{enumerate}
    \\ \hline
    Postcondición & El registro seleccionado queda actualizado o eliminado permanentemente de la base de datos, y los cambios se propagan de inmediato a las estadísticas y vistas del usuario. \\ \hline

    \caption{CU7: Gestionar Registros}
    \label{tab:esp_cu7} \\
\end{xltabular}
\endgroup


\begingroup
\small 
\renewcommand{\arraystretch}{1.4}
\begin{xltabular}{\textwidth}{|>{\columncolor{gray!20}\bfseries}p{3.5cm}|X|}
    \rowcolor{UCO}
    \textcolor{white}{Campo} & \textcolor{white}{Descripción} \\ \hline
    \endfirsthead

    \multicolumn{2}{c}{{\bfseries \dots continuación de la página anterior}} \\
    \hline
    \rowcolor{UCO}
    \textcolor{white}{Campo} & \textcolor{white}{Descripción} \\ \hline
    \endhead

    % CONTENIDO DE LA TABLA
    Nombre & Gestionar Etiquetas \\ \hline
    Identificador & CU8 \\ \hline
    Descripción & El usuario crea nuevas etiquetas personalizadas o edita/elimina las existentes dentro de su catálogo personal. \\ \hline
    Actor Principal & Usuario Registrado \\ \hline
    Actor Secundario & Ninguno \\ \hline
    Precondición & El usuario debe estar autenticado en el sistema. \\ \hline
    Flujo Principal & 
    \begin{enumerate}[leftmargin=0.5cm, topsep=0pt, partopsep=0pt, parsep=0pt, itemsep=4pt]
        \item El caso de uso comienza cuando el usuario accede a la sección de etiquetas.
        \item El sistema muestra un formulario solicitando un nombre identificativo y la selección de un color para la etiqueta.
        \item El usuario introduce el nombre deseado y selecciona un color.
        \item El usuario selecciona ''Añadir''.
        \item El sistema verifica que el nombre no esté vacío y que no exista ya una etiqueta activa con el mismo nombre para este usuario.
        \item El sistema registra la nueva etiqueta en la base de datos asignándola al identificador del usuario.
        \item El sistema actualiza la vista del catálogo mostrando la nueva etiqueta disponible para ser usada en futuros registros diarios.
    \end{enumerate} 
    \\ \hline
    Flujo Alternativo & 
    \begin{enumerate}[leftmargin=0.5cm, topsep=0pt, partopsep=0pt, parsep=0pt, itemsep=4pt]
        \item \textbf{Editar Etiqueta}: En el paso 1, el usuario pulsa ''Editar'' en una etiqueta existente. El sistema muestra el formulario con los datos actuales. El usuario los modifica y guarda. El sistema valida (paso 5) y actualiza la información en la base de datos.
        \item \textbf{Eliminar Etiqueta (Borrado Lógico)}: En el paso 1, el usuario pulsa ''Eliminar'' en una etiqueta existente. Tras solicitar confirmación, el sistema no borra el registro de la base de datos para preservar la integridad de los registros pasados y en su lugar, marca la etiqueta como inactiva, ocultándola del catálogo para que no pueda ser seleccionada en nuevos registros.
        \item \textbf{Nombre Duplicado o Inválido}: Se activa en el paso 5 si el nombre está vacío o ya existe en el catálogo activo del usuario. El sistema detiene el proceso y muestra un mensaje de error indicando que el nombre debe ser único.
        \item \textbf{Cancelar}: En cualquier momento, el usuario selecciona ''Cancelar''. El sistema cancela el proceso y regresa al catálogo principal.
    \end{enumerate}
    \\ \hline
    Postcondición & El catálogo de etiquetas del usuario queda actualizado, preservando en todo momento la integridad de los registros históricos anteriores. \\ \hline

    \caption{CU8: Gestionar Etiquetas}
    \label{tab:esp_cu8} \\
\end{xltabular}
\endgroup

\begingroup
\small 
\renewcommand{\arraystretch}{1.4}
\begin{xltabular}{\textwidth}{|>{\columncolor{gray!20}\bfseries}p{3.5cm}|X|}
    \rowcolor{UCO}
    \textcolor{white}{Campo} & \textcolor{white}{Descripción} \\ \hline
    \endfirsthead

    \multicolumn{2}{c}{{\bfseries \dots continuación de la página anterior}} \\
    \hline
    \rowcolor{UCO}
    \textcolor{white}{Campo} & \textcolor{white}{Descripción} \\ \hline
    \endhead

    % CONTENIDO DE LA TABLA
    Nombre & Visualizar Dashboard \\ \hline
    Identificador & CU9 \\ \hline
    Descripción & El usuario accede al panel principal de la aplicación para consultar un resumen de su actividad reciente, su racha actual y acceder rápidamente a la creación de registros. \\ \hline
    Actor Principal & Usuario Registrado \\ \hline
    Actor Secundario & Ninguno \\ \hline
    Precondición & El usuario debe estar autenticado en el sistema. \\ \hline
    Flujo Principal & 
    \begin{enumerate}[leftmargin=0.5cm, topsep=0pt, partopsep=0pt, parsep=0pt, itemsep=4pt]
        \item El caso de uso comienza cuando el usuario inicia sesión exitosamente o navega hacia la pantalla de inicio de la aplicación.
        \item El sistema evalúa la franja horaria local del cliente (mañana, tarde o noche) y verifica si existe un registro creado en la fecha actual.
        \item El sistema genera un mensaje de bienvenida dinámico en base a los parámetros anteriores.
        \item El sistema carga la racha actual del usuario.
        \item El sistema recupera un listado cronológico resumido con las últimas entradas del historial del usuario.
        \item El sistema renderiza el Dashboard mostrando: el mensaje de bienvenida, el contador de racha, el historial reciente y un botón destacado de acceso directo para crear un nuevo registro.
    \end{enumerate} 
    \\ \hline
    Flujo Alternativo & 
    \begin{enumerate}[leftmargin=0.5cm, topsep=0pt, partopsep=0pt, parsep=0pt, itemsep=4pt]
        \item \textbf{Usuario Sin Historial}: Si en el paso 4 y 5 el sistema detecta que el usuario no tiene ningún registro previo, el contador de racha se inicializa a 0. En la sección de entradas recientes, el sistema muestra un mensaje invitando al usuario a crear su primer registro de estado de ánimo, permitiendo extender el flujo hacia el ''CU6: Crear Registro Diario''.
    \end{enumerate}
    \\ \hline
    Postcondición & El usuario visualiza la información actualizada de su progreso y dispone de los controles necesarios para seguir interactuando con la plataforma. \\ \hline

    \caption{CU9: Visualizar Dashboard}
    \label{tab:esp_cu9} \\
\end{xltabular}
\endgroup

\begingroup
\small 
\renewcommand{\arraystretch}{1.4}
\begin{xltabular}{\textwidth}{|>{\columncolor{gray!20}\bfseries}p{3.5cm}|X|}
    \rowcolor{UCO}
    \textcolor{white}{Campo} & \textcolor{white}{Descripción} \\ \hline
    \endfirsthead

    \multicolumn{2}{c}{{\bfseries \dots continuación de la página anterior}} \\
    \hline
    \rowcolor{UCO}
    \textcolor{white}{Campo} & \textcolor{white}{Descripción} \\ \hline
    \endhead

    % CONTENIDO DE LA TABLA
    Nombre & Visualizar Historial \\ \hline
    Identificador & CU10 \\ \hline
    Descripción & El usuario consulta sus registros pasados a través de una interfaz de calendario para visualizar el detalle de cada entrada. \\ \hline
    Actor Principal & Usuario Registrado \\ \hline
    Actor Secundario & Ninguno \\ \hline
    Precondición & El usuario debe estar autenticado en el sistema. \\ \hline
    Flujo Principal & 
    \begin{enumerate}[leftmargin=0.5cm, topsep=0pt, partopsep=0pt, parsep=0pt, itemsep=4pt]
        \item El caso de uso comienza cuando el usuario selecciona el apartado ''Historial'' en el menú de navegación.
        \item El sistema renderiza una vista de calendario, estableciendo por defecto el mes y año actuales.
        \item El sistema consulta la base de datos y resalta visualmente los días que contienen un registro.
        \item El usuario selecciona un día resaltado del calendario.
        \item El sistema recupera y muestra el detalle completo del registro (puntuación global, etiquetas y comentario).
    \end{enumerate} 
    \\ \hline
    Flujo Alternativo & 
    \begin{enumerate}[leftmargin=0.5cm, topsep=0pt, partopsep=0pt, parsep=0pt, itemsep=4pt]
        \item \textbf{Navegar entre meses}: En el paso 4, el usuario utiliza los controles de ''Mes anterior'' o ''Mes siguiente''. El sistema actualiza la vista al mes seleccionado y vuelve al paso 3.
        \item \textbf{Seleccionar día sin registro}: En el paso 4, si el usuario selecciona un día que no tiene registro, el sistema muestra un mensaje indicándolo y ofrece la opción de extender el flujo hacia el ''CU6: Crear Registro Diario''.
        \item \textbf{Gestionar Registro}: En el paso 5, si el usuario decide editar o borrar la información del registro visualizado, el sistema extiende el flujo hacia el ''CU7: Gestionar Registros''.
    \end{enumerate}
    \\ \hline
    Postcondición & El usuario visualiza un calendario con el histórico de sus registros. \\ \hline
    
    \caption{CU10: Visualizar Historial}
    \label{tab:esp_cu10} \\
\end{xltabular}
\endgroup

\begingroup
\small 
\renewcommand{\arraystretch}{1.4}
\begin{xltabular}{\textwidth}{|>{\columncolor{gray!20}\bfseries}p{3.5cm}|X|}
    \rowcolor{UCO}
    \textcolor{white}{Campo} & \textcolor{white}{Descripción} \\ \hline
    \endfirsthead

    \multicolumn{2}{c}{{\bfseries \dots continuación de la página anterior}} \\
    \hline
    \rowcolor{UCO}
    \textcolor{white}{Campo} & \textcolor{white}{Descripción} \\ \hline
    \endhead

    % CONTENIDO DE LA TABLA
    Nombre & Consultar Evolución Temporal \\ \hline
    Identificador & CU11 \\ \hline
    Descripción & El usuario visualiza un gráfico temporal interactivo que muestra la progresión de su puntuación global de estado de ánimo. \\ \hline
    Actor Principal & Usuario Registrado \\ \hline
    Actor Secundario & Ninguno \\ \hline
    Precondición & El usuario debe estar autenticado en el sistema. \\ \hline
    Flujo Principal & 
    \begin{enumerate}[leftmargin=0.5cm, topsep=0pt, partopsep=0pt, parsep=0pt, itemsep=4pt]
        \item El caso de uso comienza cuando el usuario accede a la vista de ''Estadísticas''.
        \item El sistema renderiza la interfaz y establece por defecto el filtro temporal en vista ''Semanal'' (últimos 7 días).
        \item El sistema consulta la base de datos para recuperar las puntuaciones globales de los registros de ese periodo.
        \item El sistema genera y muestra un gráfico de líneas que conecta cronológicamente los datos recuperados.
    \end{enumerate} 
    \\ \hline
    Flujo Alternativo & 
    \begin{enumerate}[leftmargin=0.5cm, topsep=0pt, partopsep=0pt, parsep=0pt, itemsep=4pt]
        \item \textbf{Cambiar Filtro Temporal}: En el paso 4, el usuario selecciona un rango diferente (Mensual o Anual). El sistema actualiza el rango de fechas y vuelve al paso 3.
        \item \textbf{Sin Registros}: Se activa en el paso 3. Si no existen registros el sistema oculta el gráfico y muestra un mensaje de aviso.
    \end{enumerate}
    \\ \hline
    Postcondición & El usuario visualiza un gráfico con la evolución de su estado de ánimo. \\ \hline
    \caption{CU11: Consultar Evolución Temporal}
    \label{tab:esp_cu11} \\
\end{xltabular}
\endgroup


\begingroup
\small 
\renewcommand{\arraystretch}{1.4}
\begin{xltabular}{\textwidth}{|>{\columncolor{gray!20}\bfseries}p{3.5cm}|X|}
    \rowcolor{UCO}
    \textcolor{white}{Campo} & \textcolor{white}{Descripción} \\ \hline
    \endfirsthead

    \multicolumn{2}{c}{{\bfseries \dots continuación de la página anterior}} \\
    \hline
    \rowcolor{UCO}
    \textcolor{white}{Campo} & \textcolor{white}{Descripción} \\ \hline
    \endhead

    % CONTENIDO DE LA TABLA
    Nombre & Visualizar Resumen Estadístico \\ \hline
    Identificador & CU12 \\ \hline
    Descripción & El usuario accede a la página de análisis y visualiza distintas estadísticas sobre sus registros diarios y etiquetas. \\ \hline
    Actor Principal & Usuario Registrado \\ \hline
    Actor Secundario & Ninguno \\ \hline
    Precondición & El usuario debe estar autenticado en el sistema. \\ \hline
    Flujo Principal & 
    \begin{enumerate}[leftmargin=0.5cm, topsep=0pt, partopsep=0pt, parsep=0pt, itemsep=4pt]
        \item El caso de uso comienza cuando el usuario accede a la vista de ''Estadísticas''.
        \item El sistema procesa los registros y etiquetas del usuario para calcular tres métricas clave:
        \begin{itemize}[leftmargin=0.5cm, topsep=2pt, itemsep=2pt]
            \item Las etiquetas con mayor desviación positiva y negativa respecto a la puntuación global media de todos los registros históricos del usuario.
            \item El promedio histórico de las puntuaciones globales de los registros diarios del usuario agrupado por cada día de la semana.
            \item La distribución porcentual de los registros en distintos rangos de ánimo (días buenos, neutros y malos).
        \end{itemize}
        \item El sistema renderiza simultáneamente los componentes visuales: dos recuadros destacados para las etiquetas de mayor impacto positivo/negativo, un gráfico de barras para los días de la semana y un gráfico circular para la distribución porcentual.
        \item El usuario visualiza la información consolidada.
    \end{enumerate} 
    \\ \hline
    Flujo Alternativo & 
    \begin{enumerate}[leftmargin=0.5cm, topsep=0pt, partopsep=0pt, parsep=0pt, itemsep=4pt]
        \item \textbf{Sin Registros con Etiquetas}: Si en el paso 2 el sistema detecta que el usuario no dispone de al menos un registro etiquetas asociadas mostrará un mensaje de aviso invitando al usuario a registrar su día.
    \end{enumerate}
    \\ \hline
    Postcondición & El usuario adquiere consciencia sobre sus patrones emocionales y conductuales de forma rápida y automatizada. \\ \hline

    \caption{CU12: Visualizar Resumen Estadístico}
    \label{tab:esp_cu12} \\
\end{xltabular}
\endgroup

\begingroup
\small 
\renewcommand{\arraystretch}{1.4}
\begin{xltabular}{\textwidth}{|>{\columncolor{gray!20}\bfseries}p{3.5cm}|X|}
    \rowcolor{UCO}
    \textcolor{white}{Campo} & \textcolor{white}{Descripción} \\ \hline
    \endfirsthead

    \multicolumn{2}{c}{{\bfseries \dots continuación de la página anterior}} \\
    \hline
    \rowcolor{UCO}
    \textcolor{white}{Campo} & \textcolor{white}{Descripción} \\ \hline
    \endhead

    % CONTENIDO DE LA TABLA
    Nombre & Analizar Impacto de Etiquetas \\ \hline
    Identificador & CU13 \\ \hline
    Descripción & El usuario consulta cómo afectan sus distintas actividades (etiquetas) a su estado de ánimo general. \\ \hline
    Actor Principal & Usuario Registrado \\ \hline
    Actor Secundario & Ninguno \\ \hline
    Precondición & El usuario debe estar autenticado en el sistema y se encuentra dentro de la página de análisis. \\ \hline
    Flujo Principal & 
    \begin{enumerate}[leftmargin=0.5cm, topsep=0pt, partopsep=0pt, parsep=0pt, itemsep=4pt]
        \item El caso de uso comienza cuando el usuario accede a la sección de ''Tus Aliados''.
        \item El sistema muestra un selector con el catálogo de etiquetas del usuario.
        \item El usuario elige una etiqueta en el selector.
        \item El sistema calcula y genera un gráfico de barras que compara visualmente la puntuación global media de todos los registros del usuario frente al promedio específico de los días que contienen la etiqueta seleccionada.
        \item El sistema renderiza el gráfico y muestra un mensaje de texto interpretativo basado en la diferencia matemática entre la puntuación global media y la puntuación de la etiqueta analizada.
        \item El usuario visualiza si la etiqueta afecta positiva o negativamente en sus días.
    \end{enumerate} 
    \\ \hline
    Flujo Alternativo & 
    \begin{enumerate}[leftmargin=0.5cm, topsep=0pt, partopsep=0pt, parsep=0pt, itemsep=4pt]
        \item \textbf{Ausencia de Etiquetas}: Si en el paso 2 el sistema detecta que el usuario no tiene etiquetas registradas en su historial, deshabilita el selector y muestra un mensaje de aviso indicando que necesita vincular etiquetas a sus registros para utilizar esta función.
    \end{enumerate}
    \\ \hline
    Postcondición & El usuario obtiene una conclusión visual sobre la influencia positiva o negativa de las etiquetas en su bienestar emocional. \\ \hline

    \caption{CU13: Analizar Impacto de Etiquetas}
    \label{tab:esp_cu13} \\
\end{xltabular}
\endgroup

\begingroup
\small 
\renewcommand{\arraystretch}{1.4}
\begin{xltabular}{\textwidth}{|>{\columncolor{gray!20}\bfseries}p{3.5cm}|X|}
    \rowcolor{UCO}
    \textcolor{white}{Campo} & \textcolor{white}{Descripción} \\ \hline
    \endfirsthead

    \multicolumn{2}{c}{{\bfseries \dots continuación de la página anterior}} \\
    \hline
    \rowcolor{UCO}
    \textcolor{white}{Campo} & \textcolor{white}{Descripción} \\ \hline
    \endhead

    % CONTENIDO DE LA TABLA
    Nombre & Personalizar Escala \\ \hline
    Identificador & CU14 \\ \hline
    Descripción & El usuario personaliza los textos descriptivos de cada puntuación (del 1 al 10) y elige el conjunto de iconos visuales para sus registros. \\ \hline
    Actor Principal & Usuario Registrado \\ \hline
    Actor Secundario & Ninguno \\ \hline
    Precondición & El usuario debe estar autenticado en el sistema. \\ \hline
    Flujo Principal & 
    \begin{enumerate}[leftmargin=0.5cm, topsep=0pt, partopsep=0pt, parsep=0pt, itemsep=4pt]
        \item El caso de uso comienza cuando el usuario accede a la página de personalización y pulsa en ''Personalizar Escala de Ánimo''.
        \item El sistema muestra un formulario con los diez niveles numéricos con sus textos descriptivos actuales y un selector visual con los conjuntos de iconos predefinidos disponibles.
        \item El usuario modifica los textos descriptivos de los niveles numéricos que desee y selecciona un nuevo conjunto de iconos.
        \item El usuario selecciona la opción ''Guardar Cambios''.
        \item El sistema valida que los textos no excedan el límite máximo de caracteres permitido.
        \item El sistema actualiza la información en la base de datos y muestra un mensaje de éxito.
    \end{enumerate} 
    \\ \hline
    Flujo Alternativo & 
    \begin{enumerate}[leftmargin=0.5cm, topsep=0pt, partopsep=0pt, parsep=0pt, itemsep=4pt]
        \item \textbf{Errores de Validación}: Si en el paso 5 un texto supera la longitud máxima permitida, el sistema detiene el proceso de guardado y muestra un mensaje de aviso.
        \item \textbf{Restablecer Valores por Defecto}: En cualquier momento, el usuario puede seleccionar la opción ''Restablecer''. El sistema solicita confirmación y, tr utilizará a partir de este momento la nueva escala personalizadaas ser aceptada, restaura los textos genéricos y el conjunto de iconos original de la aplicación, guardando estos cambios en la base de datos.
    \end{enumerate}
    \\ \hline
    Postcondición & La escala del estado de animo del usuario queda actualizada en la base de datos y está disponible en la interfaz de creación y edición de registros del usuario. \\ \hline

    \caption{CU14: Personalizar Escala}
    \label{tab:esp_cu14} \\
\end{xltabular}
\endgroup

\begingroup
\small 
\renewcommand{\arraystretch}{1.4}
\begin{xltabular}{\textwidth}{|>{\columncolor{gray!20}\bfseries}p{3.5cm}|X|}
    \rowcolor{UCO}
    \textcolor{white}{Campo} & \textcolor{white}{Descripción} \\ \hline
    \endfirsthead

    \multicolumn{2}{c}{{\bfseries \dots continuación de la página anterior}} \\
    \hline
    \rowcolor{UCO}
    \textcolor{white}{Campo} & \textcolor{white}{Descripción} \\ \hline
    \endhead

    % CONTENIDO DE LA TABLA
    Nombre & Personalizar Interfaz \\ \hline
    Identificador & CU15 \\ \hline
    Descripción & El usuario ajusta la apariencia visual de la aplicación alternando entre modos de visualización (claro/oscuro) y seleccionando un color de acento principal. \\ \hline
    Actor Principal & Usuario Registrado \\ \hline
    Actor Secundario & Ninguno \\ \hline
    Precondición & El usuario debe estar autenticado en el sistema. \\ \hline
    Flujo Principal & 
    \begin{enumerate}[leftmargin=0.5cm, topsep=0pt, partopsep=0pt, parsep=0pt, itemsep=4pt]
        \item El caso de uso comienza cuando el usuario accede a la página de personalización y pulsa en ''Personalizar Interfaz'' interactúa con el panel de apariencia visual.
        \item El sistema renderiza los controles disponibles: un selector de tema (Claro, Oscuro) y una paleta con colores de acento seleccionables.
        \item El usuario selecciona una nueva opción de tema y/o un nuevo color principal.
        \item El sistema aplica el cambio visual de forma inmediata en la interfaz para que el usuario pueda previsualizar el resultado en tiempo real y actualiza las preferencias del usuario en la base de datos.
        \item El usuario visualiza los cambios en la interfaz.
    \end{enumerate} 
    \\ \hline
    Flujo Alternativo & 
    \begin{enumerate}[leftmargin=0.5cm, topsep=0pt, partopsep=0pt, parsep=0pt, itemsep=4pt]
        \item \textbf{Error Durante el Proceso}: En el paso 4 si ocurre algun problema con la actualización o guardado de la personalización del usuario el sistema mostrará un mensaje de aviso.
    \end{enumerate}
    \\ \hline
    Postcondición & La aplicación se renderiza con los nuevos parámetros estéticos seleccionados y el usuario visualiza la nueva interfaz. \\ \hline

    \caption{CU15: Personalizar Interfaz}
    \label{tab:esp_cu15} \\
\end{xltabular}
\endgroup